\chapter{Metode Penelitian}

Pada bab ini dijelaskan mengenai metode yang digunakan dalam penelitian, meliputi subjek dan objek penelitian, alat dan bahan yang digunakan, serta alur penelitian yang dilakukan secara sistematis.

\section{Subjek dan Objek Penelitian}
Dalam penelitian ini, subjek dan objek penelitian didefinisikan untuk memperjelas ruang lingkup serta fokus kajian yang dilakukan. Pada bidang sains data, subjek penelitian umumnya berupa data, algoritma, atau sistem yang digunakan dalam proses analisis dan pemodelan, bukan melibatkan partisipasi manusia secara langsung.

Subjek dalam penelitian ini adalah dataset diabetes yang digunakan sebagai sumber data utama dalam proses analisis dan pemodelan \eng{machine learning}. Dataset tersebut berisi berbagai atribut yang berkaitan dengan kondisi kesehatan individu, seperti kadar glukosa, tekanan darah, BMI, usia, serta variabel lain yang relevan dalam menentukan risiko T2DM. Selain data, subjek penelitian juga mencakup algoritma \eng{machine learning} yang digunakan, yaitu algoritma \eng{baseline} (\eng{Logistic Regression}, SVM, dan KNN) serta algoritma \eng{boosting} (\eng{Gradient Boosting}, XGBoost, dan LightGBM).

Sementara itu, objek penelitian dalam penelitian ini adalah performa model \eng{machine learning} dalam melakukan prediksi T2DM. Fokus utama penelitian adalah menganalisis dan membandingkan kinerja berbagai algoritma dalam mengklasifikasikan risiko diabetes, termasuk pengaruh \eng{hyperparameter tuning} terhadap peningkatan performa model, serta bagaimana interpretasi model dapat dilakukan menggunakan pendekatan XAI. Dengan demikian, penelitian ini berfokus pada bagaimana kombinasi algoritma, teknik optimasi, dan interpretasi model dapat menghasilkan sistem prediksi yang akurat dan dapat dipahami, khususnya dalam konteks data kesehatan.

\section{Alat dan Bahan Penelitian}

\subsection{Alat Penelitian}
Alat yang digunakan dalam penelitian ini terdiri dari perangkat keras dan perangkat lunak yang mendukung proses pengolahan data serta pemodelan \eng{machine learning}. Rinciannya adalah sebagai berikut:
\begin{enumerate}
\item \textbf{Perangkat keras}
  \begin{enumerate}[label=\alph*.]
  \item Laptop Asus Vivobook M1403QA.
  \item Prosesor AMD Ryzen 7 5800H with Radeon Graphics.
  \item Memori (RAM) 16 GB.
  \item Penyimpanan SSD 512 GB.
  \end{enumerate}
\item \textbf{Perangkat lunak}
  \begin{enumerate}[label=\alph*.]
  \item Visual Studio Code sebagai lingkungan pengembangan.
  \item Microsoft Word untuk penyusunan laporan penelitian.
  \item Microsoft Excel untuk pengolahan data awal.
  \item Mendeley sebagai manajemen referensi.
  \end{enumerate}
\end{enumerate}

\subsection{Bahan Penelitian}
Bahan penelitian yang digunakan adalah dataset EHR dari kompetisi \eng{Practice Fusion Diabetes Classification} yang diselenggarakan melalui Kaggle pada 2012 sebagai bagian dari \eng{data mining challenge} di bidang kesehatan yang berfokus pada prediksi T2DM. Dataset ini berasal dari Practice Fusion yang menyediakan data rekam medis pasien untuk keperluan analisis dan pengembangan model prediksi penyakit.

Dataset terdiri dari beberapa berkas utama, yaitu \texttt{diagnosis.csv} sebanyak 9.948 data, \texttt{medication.csv} sebanyak 9.836 data, \texttt{patient.csv} sebanyak 9.948 data, dan \texttt{transcript.csv} sebanyak 9.948 data. Berkas tersebut memuat informasi kondisi pasien, riwayat pengobatan, data demografis, spesialisasi dokter, serta catatan medis pasien. Variabel target penelitian adalah status diagnosis diabetes.

\section{Diagram Alir Penelitian}
Penelitian ini menggunakan pendekatan OSEMN (\eng{Obtain, Scrub, Explore, Model, iNterpret}) sebagai kerangka kerja dalam proses pengolahan data dan pembangunan model \eng{machine learning}. Pendekatan ini dipilih karena mampu menggambarkan alur kerja sains data secara sistematis, mulai dari pengumpulan data hingga interpretasi hasil model.

\begin{figure}[H]
\centering
\includegraphics[height=0.68\textheight]{diagram_osemn.png}
\caption{Diagram Alir Penelitian}
\label{fig:diagram-osemn}
\end{figure}

\subsection{OBTAIN (Perolehan Data)}
Tahap \eng{Obtain} merupakan tahap awal dalam kerangka kerja OSEMN yang bertujuan memperoleh data yang relevan dengan tujuan penelitian, yaitu membangun model prediksi T2DM berbasis data EHR. Pada tahap ini dilakukan identifikasi sumber data, pemahaman struktur dataset, serta penentuan variabel yang akan digunakan dalam proses pemodelan \eng{machine learning}.

Dataset yang digunakan terdiri dari beberapa berkas utama yang merepresentasikan berbagai aspek rekam medis pasien. Untuk menghasilkan dataset yang komprehensif, seluruh data diintegrasikan menggunakan atribut kunci \texttt{PatientGuid}. Ringkasan struktur dataset disajikan pada Tabel \ref{tab:struktur-dataset}.

\begin{table}[H]
\centering
\caption{Struktur Dataset}
\label{tab:struktur-dataset}
\small
\begin{tabularx}{\textwidth}{|L{3cm}|Y|Y|}
\hline
\textbf{Dataset} & \textbf{Deskripsi} & \textbf{Atribut}\\ \hline
\texttt{patient.csv} & Data demografis pasien dan variabel target. Deskripsi lengkap disajikan pada Lampiran 1. & \texttt{PatientGuid}, \texttt{Age}, \texttt{Gender}, \texttt{State}, \texttt{DMIndicator}.\\ \hline
\texttt{diagnosis.csv} & Riwayat diagnosis berdasarkan ICD-9. Deskripsi lengkap disajikan pada Lampiran 2. & \texttt{PatientGuid}, \texttt{DiagnosisCount}, \texttt{VisitCount}, \texttt{DiagnosisFreq}, \texttt{AcuteCount}, \texttt{AcuteFreq}, dan kelompok \texttt{Icd9\_001-899}.\\ \hline
\texttt{medication.csv} & Data penggunaan obat pasien. Deskripsi lengkap disajikan pada Lampiran 3. & Fitur penggunaan obat pasien yang merepresentasikan konsumsi berbagai jenis obat, dengan jumlah atribut lebih dari 2.000 kolom.\\ \hline
\texttt{transcript.csv} & Data pemeriksaan klinis. Deskripsi lengkap disajikan pada Lampiran 4. & BMI (mean, max, min), tekanan darah sistolik dan diastolik, berat badan, suhu tubuh, serta fitur perubahan.\\ \hline
\end{tabularx}
\end{table}

Berdasarkan Tabel \ref{tab:struktur-dataset}, dataset mencakup data demografis, riwayat diagnosis, penggunaan obat, dan hasil pemeriksaan klinis. Dataset \texttt{patient.csv} berfungsi sebagai sumber utama karena memuat variabel target, sedangkan dataset lainnya berperan sebagai variabel prediktor yang memperkaya informasi pasien. Seluruh dataset diintegrasikan secara bertahap menggunakan \texttt{PatientGuid} sebagai kunci penggabungan.

\subsection{SCRUB (Pembersihan dan Penyiapan Data)}
Tahap \eng{Scrub} merupakan proses pembersihan dan penyiapan data agar siap digunakan dalam analisis dan pemodelan. Dataset EHR umumnya memiliki karakteristik kompleks, seperti adanya \eng{missing values}, data duplikat, serta distribusi kelas yang tidak seimbang. Kualitas data pada tahap ini akan memengaruhi performa model yang dihasilkan.

\subsubsection{Data Cleaning}
Pembersihan data dilakukan dengan memeriksa \eng{missing values}, data duplikat, dan konsistensi nilai pada setiap fitur. Pada data EHR, nilai hilang tidak selalu merupakan kesalahan karena dapat merepresentasikan pemeriksaan yang tidak dilakukan. Nilai hilang pada fitur numerik diimputasi menggunakan median karena lebih tahan terhadap \eng{outlier} dibandingkan rata-rata. Selain itu, ditambahkan \eng{missing indicator} agar model dapat menangkap informasi bahwa suatu nilai merupakan hasil imputasi. Data duplikat diidentifikasi dan dihapus, sedangkan \eng{outlier} diperiksa berdasarkan konteks klinis dan distribusi fitur.

\subsubsection{Feature Engineering dan Encoding}
Variabel kategorikal dikonversi menjadi numerik menggunakan \eng{one-hot encoding}, khususnya pada fitur nominal seperti jenis kelamin dan kategori diagnosis. Tipe data setiap fitur juga disesuaikan dengan kebutuhan model.

Fitur numerik kontinu distandarisasi menggunakan \eng{StandardScaler} untuk menyamakan skala antarvariabel. Proses ini diterapkan pada usia, BMI, tekanan darah, suhu tubuh, jumlah diagnosis, dan jumlah kunjungan. Fitur biner seperti kelompok ICD dan indikator penggunaan obat tidak distandarisasi. Bagi algoritma berbasis pohon, standarisasi bukan kebutuhan utama, tetapi tetap dimasukkan dalam \eng{pipeline} yang sesuai agar perbandingan antarmodel tetap konsisten.

\subsubsection{Pembagian Data}
Dataset dibagi menjadi data latih dan data uji dengan metode \eng{stratified split} pada proporsi 70:30. Data latih digunakan untuk pelatihan dan \eng{hyperparameter tuning}, sedangkan data uji digunakan secara eksklusif untuk evaluasi akhir. Pemisahan ini dilakukan untuk mencegah \eng{data leakage} dan menilai kemampuan generalisasi model terhadap data baru.

\subsubsection{Penanganan Class Imbalance}
Distribusi kelas pada dataset kesehatan sering tidak seimbang karena jumlah data non-diabetes dapat lebih besar dibandingkan data T2DM. Pendekatan yang digunakan meliputi penyesuaian \eng{class weight} pada model dan teknik \eng{oversampling} seperti SMOTE. Pemilihan metode dilakukan berdasarkan hasil evaluasi agar model tidak hanya akurat secara keseluruhan, tetapi juga mampu mendeteksi kelas diabetes dengan baik.

\subsection{EXPLORE (Eksplorasi Data)}
Tahap \eng{Explore} bertujuan memahami karakteristik data melalui \eng{Exploratory Data Analysis} (EDA) sebelum pemodelan. Analisis dilakukan menggunakan statistik deskriptif dan visualisasi untuk mengidentifikasi distribusi data, pola hubungan antarvariabel, serta indikasi awal faktor-faktor yang berpengaruh terhadap status T2DM.

\subsubsection{Analisis Distribusi Kelas}
Analisis distribusi kelas dilakukan dengan menghitung jumlah dan persentase kelas diabetes dan non-diabetes serta mengidentifikasi adanya \eng{class imbalance}. Hasil analisis digunakan untuk menentukan strategi penanganan ketidakseimbangan data dan metrik evaluasi yang sesuai.

\subsubsection{Statistik Deskriptif Variabel}
Statistik deskriptif digunakan untuk menggambarkan nilai rata-rata, minimum, maksimum, standar deviasi, dan distribusi fitur utama seperti BMI, tekanan darah, dan indikator klinis lainnya. Histogram dan \eng{boxplot} digunakan untuk memahami sebaran data serta mendeteksi potensi anomali atau \eng{outlier}.

\subsubsection{Analisis Hubungan Antarvariabel}
Analisis dilakukan untuk mengidentifikasi hubungan antara variabel prediktor dan status diabetes serta hubungan antarf fitur. Pendekatan yang digunakan bersifat deskriptif, misalnya membandingkan distribusi fitur pada kelompok diabetes dan non-diabetes dan menggunakan \eng{heatmap} korelasi. Analisis ini tidak digunakan untuk menarik kesimpulan kausal.

\subsection{MODEL (Pemodelan)}
Tahap \eng{Model} merupakan tahap inti yang bertujuan membangun model klasifikasi untuk memprediksi status T2DM berdasarkan data EHR. Dua kelompok algoritma digunakan, yaitu algoritma \eng{baseline} dan algoritma \eng{boosting}. Algoritma \eng{baseline} terdiri atas \eng{Logistic Regression}, SVM, dan KNN, sedangkan algoritma \eng{boosting} terdiri atas \eng{Gradient Boosting}, XGBoost, dan LightGBM.

Penggunaan dua kelompok tersebut bertujuan melihat peningkatan performa yang dihasilkan pendekatan \eng{ensemble} dibandingkan metode dasar. Seluruh pelatihan dan optimasi dilakukan hanya pada data latih, sedangkan data uji digunakan untuk evaluasi akhir.

\subsubsection{Desain Eksperimen}
Enam algoritma dilatih menggunakan dataset, variabel prediktor, dan pembagian data yang sama. Metrik evaluasi yang seragam digunakan agar perbandingan performa bersifat adil. Variabel prediktor berasal dari data demografis, kondisi klinis, riwayat diagnosis, penggunaan obat, dan indikator kesehatan lainnya, sedangkan target adalah status diabetes.

\subsubsection{Pipeline Pemodelan}
Gambar \ref{fig:pipeline-ml} menunjukkan alur proses pembangunan model \eng{machine learning} untuk prediksi Diabetes Mellitus Tipe 2.

\begin{figure}[H]
\centering
\includegraphics[width=0.98\textwidth]{pipeline_ml.png}
\caption{\eng{Machine Learning Pipeline} untuk Prediksi Diabetes}
\label{fig:pipeline-ml}
\end{figure}

Proses dimulai dari \eng{data preprocessing}, dilanjutkan dengan \eng{feature engineering}, pembentukan variabel prediktor $X$ dan target $Y$, serta pembagian data secara \eng{stratified}. Optuna digunakan untuk menemukan kombinasi \eng{hyperparameter} terbaik. Model kemudian dilatih dan dievaluasi menggunakan \eng{accuracy}, \eng{precision}, \eng{recall}, \eng{F1-score}, ROC-AUC, dan \eng{cross-validation}. Model terbaik selanjutnya dianalisis menggunakan SHAP.

\subsubsection{Baseline Model}
Model \eng{baseline} digunakan untuk memperoleh gambaran performa awal. \eng{Logistic Regression} dipilih karena sederhana dan memiliki interpretasi yang baik, SVM dipilih karena kemampuannya menangani data berdimensi tinggi, sedangkan KNN digunakan sebagai metode berbasis jarak yang sederhana.

\subsubsection{Model Boosting}
Model \eng{boosting} yang digunakan adalah \eng{Gradient Boosting}, XGBoost, dan LightGBM. \eng{Gradient Boosting} meminimalkan fungsi \eng{loss} secara iteratif, XGBoost menambahkan regularisasi dan optimasi komputasi, sedangkan LightGBM menggunakan pertumbuhan pohon \eng{leaf-wise} untuk mempercepat pelatihan pada dataset besar.

\subsubsection{Tuning Hyperparameter}
Optuna digunakan untuk mengoptimalkan parameter setiap algoritma. Parameter yang ditelusuri antara lain:
\begin{itemize}
\item \eng{Logistic Regression}: parameter regularisasi.
\item SVM: jenis \eng{kernel}, parameter $C$, dan parameter kernel.
\item KNN: jumlah tetangga $k$, bobot, dan metrik jarak.
\item \eng{Gradient Boosting}: \eng{learning rate}, jumlah estimator, dan kedalaman pohon.
\item XGBoost dan LightGBM: \eng{learning rate}, \eng{max depth}, jumlah pohon, subsampling, dan parameter regularisasi.
\end{itemize}

\subsubsection{Model Final dan Evaluasi}
Model final diperoleh dari kombinasi parameter terbaik pada masing-masing algoritma. Evaluasi dilakukan pada data uji menggunakan \eng{accuracy}, \eng{precision}, \eng{recall}, \eng{F1-score}, dan ROC-AUC. \eng{Cross-validation} digunakan pada tahap pelatihan untuk menilai kestabilan model dan mengidentifikasi potensi \eng{overfitting}. SHAP hanya diterapkan pada model dengan performa terbaik.

\subsection{INTERPRET (Evaluasi dan Interpretasi Model)}
Tahap \eng{iNterpret} merupakan tahap akhir yang bertujuan mengevaluasi performa model dan memahami hasil prediksi. Evaluasi dilakukan untuk mengukur kemampuan model mengklasifikasikan status T2DM, khususnya pada data yang berpotensi tidak seimbang. Interpretasi model digunakan agar hasil prediksi dapat dipahami dalam konteks kesehatan.

\subsubsection{Evaluasi Model}
Evaluasi model dilakukan pada data uji yang belum pernah dilihat selama pelatihan.

\paragraph{Confusion Matrix.}
Komponen yang dianalisis adalah:
\begin{itemize}
\item \eng{True Positive} (TP): data diabetes yang diprediksi benar sebagai diabetes.
\item \eng{True Negative} (TN): data non-diabetes yang diprediksi benar sebagai non-diabetes.
\item \eng{False Positive} (FP): data non-diabetes yang diprediksi sebagai diabetes.
\item \eng{False Negative} (FN): data diabetes yang diprediksi sebagai non-diabetes.
\end{itemize}

\paragraph{Metrik Evaluasi.}
\eng{Accuracy} digunakan untuk mengukur ketepatan keseluruhan, \eng{precision} untuk mengukur ketepatan prediksi kelas diabetes, \eng{recall} untuk mengukur kemampuan mendeteksi kasus diabetes, dan \eng{F1-score} untuk menyeimbangkan \eng{precision} dan \eng{recall}. Dalam konteks penelitian ini, \eng{recall} menjadi metrik penting karena berkaitan dengan minimisasi \eng{false negative}.

\paragraph{ROC-AUC.}
ROC-AUC digunakan untuk menilai kemampuan model membedakan kelas diabetes dan non-diabetes pada berbagai nilai ambang. \eng{Cross-validation} digunakan untuk memastikan hasil evaluasi tidak bergantung pada satu pembagian data.

\subsubsection{Interpretasi Model}
Metode SHAP diterapkan pada model dengan performa terbaik. Visualisasi yang digunakan meliputi:
\begin{itemize}
\item \eng{Summary plot} untuk menunjukkan fitur paling berpengaruh secara global.
\item \eng{Dependence plot} untuk menunjukkan hubungan nilai fitur dan pengaruhnya terhadap prediksi.
\item \eng{Force plot} untuk menjelaskan kontribusi fitur pada prediksi individu.
\end{itemize}

Keluaran tahap ini berupa perbandingan performa seluruh algoritma, penentuan model terbaik berdasarkan metrik evaluasi, serta visualisasi dan interpretasi kontribusi fitur menggunakan SHAP. Hasil tersebut digunakan untuk menarik kesimpulan mengenai model terbaik dan memahami faktor-faktor yang memengaruhi prediksi T2DM berdasarkan data EHR.
